\documentclass[11pt,a4paper]{article}
\newcommand{\resultsdir}{../../results/task3}
% Common preamble of the task reports (compile with XeLaTeX).
\usepackage{fontspec}
\usepackage{polyglossia}
\setmainlanguage{polish}
\usepackage{amsmath,amssymb}
\usepackage[a4paper,margin=2.3cm]{geometry}
\usepackage{graphicx}
\usepackage{booktabs}
\usepackage{longtable}
\usepackage{array}
\usepackage{multirow}
\usepackage{adjustbox}
\usepackage{makecell}
\usepackage{placeins}
\usepackage{float}
\usepackage[font=small,labelfont=bf]{caption}
\usepackage{subcaption}
\usepackage{xcolor}
\usepackage{enumitem}
\usepackage[hidelinks]{hyperref}

% Without the texlive-lang-polish package XeLaTeX would hyphenate Polish
% words with English patterns, so hyphenation is switched off and looser
% inter-word spacing is allowed instead. With Polish patterns installed the
% next two penalties can be removed.
\hyphenpenalty=10000
\exhyphenpenalty=10000
\tolerance=3000
\emergencystretch=3em
\setlist{itemsep=2pt,topsep=4pt}
\setlength{\parskip}{3pt}
\renewcommand{\arraystretch}{1.12}

\newcommand{\tab}[1]{\input{\resultsdir/tables/#1}}
% table scaled down to the text width only when it is wider
\newcommand{\fittab}[1]{\adjustbox{max width=\textwidth}{\tab{#1}}}
\newcommand{\fig}[1]{\resultsdir/figures/#1}
\newcommand{\code}[1]{\texttt{#1}}
% left-aligned paragraph column (no stretched spaces in narrow cells)
\newcolumntype{L}[1]{>{\raggedright\arraybackslash}p{#1}}

% \maketitlepage{task title}{subtitle}
\newcommand{\maketitlepage}[2]{%
  \begin{titlepage}
    \centering
    {\large Politechnika Łódzka\\Wydział Fizyki Technicznej, Informatyki i Matematyki Stosowanej\par}
    \vspace{3.2cm}
    {\large Analiza danych złożonych z detekcją wyjątków\par}
    \vspace{0.8cm}
    {\LARGE\bfseries #1\par}
    \vspace{0.4cm}
    {\large #2\par}
    \vfill
    {\large Jędrzej Kostyk\par}
    {nr indeksu 264339\par}
    \vspace{1.2cm}
    {Rok akademicki 2026/2027\par}
  \end{titlepage}
}


\begin{document}
\maketitlepage{Zadanie 3}{Zespołowa detekcja wyjątków:\\projekt własnego rozwiązania}
\tableofcontents
\newpage

\section{Cel i zakres}

Celem zadania jest zaprojektowanie i ocena zespołowego systemu detekcji
wyjątków, który łączy wyniki kilku detektorów o różnych zasadach działania.
Sprawdzamy, czy agregacja przewyższa najlepszy pojedynczy detektor.
Porównano trzy grupy rozwiązań:
\begin{enumerate}
  \item \textbf{modele pojedyncze} -- cztery detektory bazowe oceniane
  niezależnie (punkt odniesienia),
  \item \textbf{głosowanie większościowe} (\emph{Majority Voting}) z progami
  od 1/4 do 4/4 głosów,
  \item \textbf{fuzja wyników} $S(x)=\sum_k w_k\,\hat s_k(x)$ z różnymi
  sposobami normalizacji, wag i wyboru detektorów, w tym dwa własne warianty:
  DSF i CBF.
\end{enumerate}
Ostatnia część raportu (sekcja~\ref{sec:synthesis}) zbiera wyniki
wszystkich trzech zadań i odpowiada na pytanie przewodnie projektu: co miało
największy wpływ na jakość detekcji -- wybór algorytmu, reprezentacja danych
czy połączenie detektorów.

Kod znajduje się w katalogu \code{task3\_ensemble/}: normalizacja, wagi
i agregacja w \code{ensemble.py}, detektory bazowe w \code{run\_base.py},
warianty zespołów w \code{run\_ensembles.py}, a tabele i rysunki
w \code{analyze.py}.

\section{Dane i protokół}

Wykorzystano te same cztery zbiory benchmarkowe co w zadaniu~2: Cardio
(ADBench~\cite{adbench}; UCI Cardiotocography~\cite{cardio}), Arrhythmia
(ODDS~\cite{odds}; UCI Arrhythmia~\cite{arrhythmia}), WDBC (ADBench; UCI
Breast Cancer Wisconsin (Diagnostic)~\cite{wdbc}) i Campaign (ADBench; UCI
Bank Marketing~\cite{campaign}). Dzięki temu porównanie z zadaniem~2 jest
bezpośrednie. Dane i \textbf{etykiety anomalii pochodzą z benchmarków}
bez żadnych modyfikacji (opis -- raporty z zadań 1 i 2).

Protokół jest taki sam we wszystkich zadaniach: te same warstwowe podziały
50/25/25\,\% (10 podziałów dla Cardio i Arrhythmii, 30 dla WDBC, 5 dla
Campaign), standaryzacja i reprezentacje dopasowywane na części uczącej,
detektory uczone bez etykiet. \textbf{Część walidacyjna} służy do doboru
progów poszczególnych detektorów, wag, podzbioru detektorów i progu decyzji
zespołu. \textbf{Część testowa} nie jest używana do żadnej decyzji.
Liczba anomalii w części walidacyjnej bardzo się różni między zbiorami:
2--3 w WDBC, ok. 17 w Arrhythmii, 44 w Cardio i 1160 w Campaign. Jak
pokażą wyniki, ma to kluczowe znaczenie dla zespołów.

Wszystkie warianty zespołów, w tym oba warianty własne, zostały
zdefiniowane z ustalonymi parametrami przed uruchomieniem ich na tych
danych. Dodatkowo oceniono je na \textbf{nowych podziałach} (kolejny blok
ziaren: 10--19 dla Cardio i Arrhythmii, 30--59 dla WDBC, 5--9 dla Campaign),
nieużywanych w żadnym etapie wyboru konfiguracji. Stanowią one test
potwierdzający (sekcja~\ref{sec:confirm}).

\section{Detektory bazowe}

Wybrano cztery detektory o różnych zasadach działania:
\begin{itemize}
  \item \textbf{Isolation Forest} -- izolacja losowymi cięciami,
  \item \textbf{LOF} -- lokalna gęstość względem sąsiadów,
  \item \textbf{One-Class SVM} -- granica obszaru danych w przestrzeni jądra RBF,
  \item \textbf{błąd rekonstrukcji PCA} -- odległość od podprzestrzeni
  głównych składowych (wykrywa odchylenia w kierunkach o małej wariancji).
\end{itemize}
Zgodnie z zaleceniem, konfiguracje wybrano na podstawie wcześniejszych
eksperymentów. Dla każdego zbioru i detektora przyjęto reprezentację
i hiperparametry o najwyższej średniej PR-AUC na walidacji w zadaniu~2
(tabela~\ref{tab:base}). Detektor bazowy to więc para
\emph{reprezentacja + algorytm}, np. OCSVM na dwuwymiarowym UMAP dla Cardio,
LOF i OCSVM na jednej składowej PCA dla WDBC, LOF i OCSVM w przestrzeni
latentnej autoenkodera dla Arrhythmii.

\begin{table}[H]
  \centering\small
  \caption{Konfiguracje detektorów bazowych (wybór na walidacji
  w zadaniu~2).}
  \label{tab:base}
  \resizebox{\textwidth}{!}{\tab{base_configs}}
\end{table}

\section{Metody agregacji}

\subsection{Ujednolicenie wyników anomalii}

Wszystkie detektory zwracają wynik, w którym wyższa wartość oznacza
większą anomalność: IForest $-$\code{score\_samples}, LOF i OCSVM
$-$\code{score\_samples}/$-$\code{decision\_function}, PCA -- kwadrat błędu
rekonstrukcji. Skale są jednak zupełnie różne (np. LOF $\approx 1$--3,
błąd rekonstrukcji od 0 do kilkuset), więc przed agregacją każdy wynik jest
normalizowany. Zbadano cztery normalizacje, każdą dopasowywaną osobno dla
każdego detektora:
\begin{itemize}
  \item \textbf{min-max}: $(s-\min)/(\max-\min)$,
  \item \textbf{z-score}: $(s-\mu)/\sigma$,
  \item \textbf{odporny z-score}: $(s-\mathrm{med})/(1{,}4826\cdot\mathrm{MAD})$
  -- mediana i MAD nie zależą od kilku skrajnych wartości, czyli od samych
  anomalii,
  \item \textbf{rangi (dystrybuanta empiryczna)}: odsetek wyników
  referencyjnych mniejszych od $s$.
\end{itemize}

Ważne jest, \textbf{na czym} dopasować normalizację. Naturalnym wyborem są
wyniki części uczącej, ale są to wyniki \emph{in-sample}: detektor ocenia
punkty, na których był uczony. Rysunek~\ref{fig:norm} pokazuje skutek dla
OCSVM z wąskim jądrem ($\gamma=10/d$) na Campaign. Punkty uczące są
wektorami nośnymi lub leżą tuż przy nich i dostają niemal identyczne, niskie
wyniki, a 71\,\% punktów testowych ma wynik wyższy niż \emph{każdy} punkt
uczący. Normalizacja rangowa dopasowana na części uczącej przypisuje im
wszystkim wartość 1 i informacja o ich kolejności ginie. Dlatego
normalizacje dopasowywano na \textbf{wynikach części walidacyjnej}
(\emph{out-of-sample}, bez użycia etykiet). Tabela~\ref{tab:norm} pokazuje,
że dla fuzji przez maksimum zmiana ta podnosi średnią PR-AUC przy rangach
z 0,412 do 0,524, a przy z-score z 0,483 do 0,576.

\begin{figure}[H]
  \centering
  \includegraphics[width=\textwidth]{\fig{normalisation}}
  \caption{OCSVM ($\gamma=10/d$) na Campaign, podział 0: rozkład surowych
  wyników części uczącej (in-sample), walidacyjnej i testowej oraz wyniki
  testowe po normalizacji rangowej dopasowanej na części uczącej (nasycenie
  w 1), rangowej dopasowanej na walidacji i odpornym z-score dopasowanym na
  walidacji.}
  \label{fig:norm}
\end{figure}

\begin{table}[H]
  \centering\small
  \caption{Wpływ normalizacji i danych, na których ją dopasowano (średnia
  PR-AUC z 4 zbiorów, część testowa): fuzja przez średnią, fuzja przez
  maksimum oraz wariant CBF (PR-AUC i F1).}
  \label{tab:norm}
  \resizebox{\textwidth}{!}{\tab{normalisation}}
\end{table}

\subsection{Majority Voting}

Każdy detektor podejmuje decyzję binarną na podstawie własnego progu,
wybranego jako maksymalizujący F1 na walidacji (wariant główny) albo jako
kwantyl $1-c$ wyników walidacyjnych, gdzie $c$ to udział anomalii w części
walidacyjnej (wariant ,,c''; z etykiet wykorzystywana jest tylko ta jedna
liczba). Obserwacja jest anomalią, gdy wskaże ją co najmniej $m$ z 4
detektorów, $m\in\{1,2,3,4\}$. Próg $m=1$ to suma (OR), a $m=4$ --
iloczyn (AND) decyzji. Ranking potrzebny do ROC-AUC i PR-AUC tworzy liczba
głosów. Jest on taki sam dla każdego $m$, więc te metryki nie zależą od progu.

\subsection{Weighted Score Fusion}

$S(x)=\sum_k w_k\,\hat s_k(x)$, $\sum_k w_k=1$, gdzie $\hat s_k$ jest
znormalizowanym wynikiem. Zbadano schematy wag wyznaczanych na walidacji:
\begin{itemize}
  \item \textbf{równe} -- zwykła średnia,
  \item \textbf{PR-AUC} -- $w_k\propto \mathrm{AP}_k^{wal}$,
  \item \textbf{skill} -- $w_k\propto \max(0,\ \mathrm{AP}_k^{wal}-\pi)$,
  gdzie $\pi$ to udział anomalii: wagę dostaje tylko przewaga nad
  losowym rankingiem, a detektor nie lepszy od losowego dostaje 0,
  \item \textbf{skill $\times$ różnorodność} -- skill pomnożony przez
  $1-\bar\rho_k$, gdzie $\bar\rho_k$ to średnia korelacja Spearmana wyników
  detektora $k$ z pozostałymi (na walidacji),
\end{itemize}
oraz \textbf{maksimum} znormalizowanych wyników (klasyczna alternatywa
dla średniej). Każdy schemat zbadano dla wszystkich detektorów i dla
podzbioru wybranego zachłannie na walidacji. Próg decyzji wyznaczano
dwoma kryteriami: (1) maksymalizacja F1 na walidacji (F1 równoważy
Precision i Recall, a próg wybrany wprost pod raportowaną metrykę jest
najprostszy do uzasadnienia), (2) kwantyl $1-c$ zfuzjowanych wyników
walidacyjnych, gdzie $c$ to udział anomalii w walidacji (nie wymaga
wskazania, które punkty są anomaliami).

\subsection{Warianty własne}
\label{sec:own}

Oba warianty własne odpowiadają na to samo pytanie -- \emph{jak dobrać
detektory i wagi, gdy jakość detektorów da się ocenić tylko na małej
walidacji} -- ale w przeciwny sposób.

\paragraph{DSF -- \emph{diversity-weighted skill fusion}.} Założenie: zespół
zyskuje na detektorach, które są dobre \emph{i} wnoszą informację
niedostępną dla innych. DSF składa się z rangowej normalizacji (dopasowanej
na walidacji), wag $w_k\propto \max(0,\mathrm{AP}_k-\pi)\cdot(1-\bar\rho_k)$
oraz zachłannego wyboru detektorów. Wybór zaczyna od detektora o najlepszej
walidacyjnej PR-AUC i dodaje kolejne, dopóki poprawiają walidacyjną PR-AUC
fuzji o więcej niż $10^{-3}$. Próg wybierany jest przez maksymalizację F1
na walidacji. DSF w pełni ufa walidacji: dopasowuje do niej zarówno skład
zespołu, jak i ciągłe wagi.

\paragraph{CBF -- \emph{competence-based fusion}.} Założenie: przy małej
walidacji ciągłe wagi i zachłanny wybór dopasowują się do szumu, a człon
różnorodności nagradza detektory słabe (słaby detektor jest z definicji mało
skorelowany z dobrymi). CBF:
\begin{enumerate}
  \item Dla każdego detektora szacuje \emph{niepewność} jego walidacyjnej
  PR-AUC metodą bootstrap (500 prób losowanych ze zwracaniem osobno z anomalii
  i obserwacji normalnych części walidacyjnej).
  \item Za \emph{kompetentne} uznaje detektory, które w co najmniej 10\,\%
  prób bootstrapowych są nie gorsze od detektora najlepszego na walidacji,
  czyli te, o których dane nie pozwalają powiedzieć, że są gorsze.
  \item Detektory kompetentne uśrednia z \emph{równymi} wagami. Równe
  wagi nie wymagają estymacji, a różnic w jakości i tak nie da się
  wiarygodnie zmierzyć. Detektory wyraźnie gorsze odrzuca.
  \item Normalizuje odpornym z-score dopasowanym na wynikach walidacyjnych,
  który nie nasyca się i nie zależy od skrajnych wartości.
  \item Próg decyzyjny wybiera przez maksymalizację F1 na walidacji.
\end{enumerate}
Gdy jeden detektor wyraźnie dominuje, CBF sprowadza się do niego. Gdy kilka
jest nieodróżnialnych, uśrednia je i zmniejsza wariancję. Jest to zasada
,,uśredniaj tylko to, czego nie umiesz rozróżnić''.

\section{Wyniki}

\subsection{Pełna tabela wyników}

\begin{center}
\captionof{table}{Wyniki wszystkich wariantów (część testowa, podziały
projektowe, średnia $\pm$ odchylenie std.). Czas -- łączny czas
przygotowania reprezentacji, uczenia i oceny wykorzystanych detektorów oraz
agregacji. Pogrubienie -- najlepsza wartość na zbiorze.}
\label{tab:main}
\end{center}
{\footnotesize\setlength{\tabcolsep}{3pt}\tab{main_results}}

Najważniejsze obserwacje z tabeli~\ref{tab:main}:
\begin{itemize}
  \item \textbf{Zespół przewyższa najlepszy pojedynczy detektor (wybrany na
  walidacji) na trzech zbiorach.} Na Cardio DSF osiąga PR-AUC 0,872, a fuzja
  z wagami skill 0,870, wobec 0,836 dla OCSVM na UMAP. Na Arrhythmii CBF
  i DSF mają 0,604 i 0,602 wobec 0,590, a na Campaign 0,320 i 0,328 wobec
  0,303. Przewagi są niewielkie (0,01--0,04), ale dotyczą zbiorów, na których
  walidacja zawiera kilkadziesiąt lub więcej anomalii.
  \item \textbf{Na WDBC wszystkie zespoły przegrywają}, mimo że ich składniki
  są dobre: fuzja średnia 0,785, CBF 0,851 wobec 0,886 najlepszego
  pojedynczego. LOF i OCSVM na jednej składowej PCA mają osobno 0,92--0,93,
  a dołączenie słabszych IForest (0,75) i rekonstrukcji (0,46) rozmywa ich
  wynik. Walidacja z 2--3 anomaliami nie pozwala wiarygodnie odrzucić
  słabszych detektorów. Nawet wybór najlepszego pojedynczego detektora jest
  tu niestabilny: w 12 z 30 podziałów walidacja wskazała IForest zamiast LOF.
  \item \textbf{Majority Voting} ma niską PR-AUC (ranking z liczby głosów ma
  tylko 5 poziomów), ale w F1 bywa konkurencyjny: MV 2/4 ma najwyższe F1 na
  Campaign (0,375 wobec 0,318 najlepszego pojedynczego), a MV 3/4 na
  Arrhythmii dorównuje najlepszemu pojedynczemu (0,532).
  \item Fuzja przez \textbf{maksimum} (z rangami) jest najgorszym schematem
  fuzji poza Campaign (średnio 0,524): wzmacnia fałszywe alarmy każdego
  detektora.
\end{itemize}

\subsection{Majority Voting: wpływ progu}

Rysunek~\ref{fig:voting} pokazuje klasyczny kompromis. Zwiększanie liczby
wymaganych głosów podnosi Precision (Cardio: 0,38 przy 1/4, 0,59 przy 2/4,
0,77 przy 3/4, 0,995 przy 4/4) i obniża Recall (0,95 $\to$ 0,82 $\to$ 0,66
$\to$ 0,35). F1 jest najwyższe przy $m=2$ lub $m=3$. Reguła AND ($m=4$) jest
zbyt restrykcyjna: na Campaign, gdzie detektory rzadko się zgadzają (Jaccard
0,07--0,46), Recall spada do 0,06. Reguła OR ($m=1$) daje najwyższy Recall,
ale sumuje fałszywe alarmy wszystkich detektorów. Głosowanie przekreśla
unikalną wiedzę jednego detektora: na Cardio OCSVM na UMAP wykrywa 75\,\%
anomalii lokalnych, a MV 3/4 tylko 54\,\%, bo pozostałe detektory ich nie
widzą. Progi wyznaczane z udziału anomalii (linie przerywane) dają przy
$m=2$ nieco wyższe F1 niż progi z walidacyjnego F1 na trzech zbiorach
(np. WDBC 0,74 wobec 0,71) -- przy bardzo małej walidacji próg z jednego
kwantyla jest stabilniejszy niż próg optymalizowany na kilku anomaliach.

\begin{figure}[H]
  \centering
  \includegraphics[width=\textwidth]{\fig{voting_threshold}}
  \caption{Majority Voting: Precision, Recall i F1 w zależności od
  wymaganej liczby głosów (linie ciągłe -- progi detektorów z walidacyjnego
  F1, przerywane -- z udziału anomalii w walidacji); linia kropkowana -- F1
  najlepszego pojedynczego detektora.}
  \label{fig:voting}
\end{figure}

\subsection{Zgodność i komplementarność detektorów}

\begin{figure}[H]
  \centering
  \includegraphics[width=0.95\textwidth]{\fig{jaccard_heatmap}}
  \caption{Indeks Jaccarda zbiorów obserwacji testowych oznaczonych przez
  pary detektorów bazowych (średnia po podziałach).}
  \label{fig:jac}
\end{figure}

Rysunek~\ref{fig:jac} pokazuje bardzo różną strukturę zgodności na
poszczególnych zbiorach:
\begin{itemize}
  \item \textbf{WDBC}: LOF i OCSVM działają na tej samej, jednej składowej PCA
  i są praktycznie identyczne (Jaccard 0,95), więc nie są komplementarne.
  IForest i rekonstrukcja PCA są od nich bardziej niezależne (0,28--0,57),
  ale słabsze. Zespół nie ma tu z czego zyskać -- dodanie słabszych,
  niezależnych detektorów tylko szkodzi.
  \item \textbf{Arrhythmia}: IForest, LOF i OCSVM są podobnej jakości
  (PR-AUC 0,57--0,60) i umiarkowanie zgodne (0,61--0,66). Uśrednienie
  zmniejsza wariancję -- CBF w połowie podziałów uśrednia wszystkie trzy.
  \item \textbf{Campaign}: detektory prawie się nie zgadzają (IForest--LOF
  0,07), a wszystkie są słabe (0,20--0,30). Niska zgodność oznacza tu
  w dużej mierze losowość błędów, ale także różne podzbiory wykrywanych
  anomalii: 24\,\% anomalii wskazuje dokładnie jeden detektor.
  \item \textbf{Cardio}: OCSVM na UMAP wyraźnie dominuje (0,84 wobec
  0,47--0,59), a jego decyzje mało pokrywają się z pozostałymi (0,32--0,42),
  bo działa w innej reprezentacji. Pozostałe detektory, choć słabsze, niosą
  informację, której OCSVM nie ma -- stąd zysk fuzji ważonej.
\end{itemize}

Rysunek~\ref{fig:comp} zestawia miary komplementarności z zyskiem zespołu
na poziomie pojedynczych podziałów. Korelacje są słabe (|$r$| $\le$ 0,38):
zysk w pojedynczym podziale jest silnie zaszumiony, szczególnie na WDBC.
Najwyraźniejszy jest związek z udziałem anomalii wykrytych przez dokładnie
jeden detektor ($r=0{,}38$). Gdy detektory wykrywają różne anomalie,
uśrednianie ma co łączyć.

\begin{figure}[H]
  \centering
  \includegraphics[width=\textwidth]{\fig{complementarity_gain}}
  \caption{Zysk zespołu względem najlepszego pojedynczego detektora (wybranego
  na walidacji) w poszczególnych podziałach a miary komplementarności:
  średni Jaccard decyzji, udział anomalii wykrytych przez dokładnie jeden
  detektor, udział fałszywych alarmów zgłoszonych przez jeden detektor
  oraz przewaga najlepszego detektora na walidacji.}
  \label{fig:comp}
\end{figure}

\subsection{Wpływ normalizacji, wag, wyboru detektorów i progu}

Tabela~\ref{tab:abl} i rys.~\ref{fig:abl} przedstawiają pełną ablację
fuzji: 4 normalizacje $\times$ 4 schematy wag $\times$ 2 sposoby wyboru
detektorów $\times$ 2 kryteria progu.

\begin{table}[H]
  \centering\small
  \caption{Ablacja fuzji wyników (normalizacja dopasowana na walidacji):
  średnia PR-AUC i F1 z 4 zbiorów (test) oraz różnica względem najlepszego
  pojedynczego detektora ($\Delta$); F1 dla progu z walidacyjnego F1 i dla
  progu z udziału anomalii ($c$); ,,Wygrane'' -- liczba zbiorów, na których
  fuzja ma wyższą PR-AUC niż najlepszy pojedynczy detektor.}
  \label{tab:abl}
  \resizebox{\textwidth}{!}{\tab{ablation}}
\end{table}

\begin{figure}[H]
  \centering
  \includegraphics[width=0.85\textwidth]{\fig{ablation_heatmap}}
  \caption{Zmiana średniej PR-AUC fuzji względem najlepszego pojedynczego
  detektora w zależności od normalizacji i wag (lewo -- wszystkie detektory,
  prawo -- zachłanny wybór).}
  \label{fig:abl}
\end{figure}

\begin{itemize}
  \item \textbf{Wybór detektorów} ma największy wpływ. Fuzja
  \emph{wszystkich} detektorów jest zwykle gorsza od najlepszego
  pojedynczego (do $-0{,}068$), a zachłanny wybór daje zysk przy min-max,
  z-score i odpornym z-score (+0,010 do +0,021) i wygrywa na 2--3 z 4 zbiorów.
  Przyczyną jest słaby detektor -- rekonstrukcja PCA (0,26--0,47) -- który
  w średniej wszystkich detektorów obniża wynik, a zachłanny wybór go pomija.
  \item \textbf{Normalizacja}: przy wyborze zachłannym min-max jest najlepszy
  (+0,016 do +0,021), a rangi najgorsze. Przy fuzji wszystkich detektorów
  odporny z-score wypada najgorzej ($-0{,}042$ do $-0{,}068$): nie ogranicza
  wartości z góry, więc pojedyncze ekstremalne wyniki słabego detektora
  dominują w średniej. Min-max i rangi ograniczają wkład każdego detektora
  do $[0,1]$. Rangi przy kilkuset punktach walidacyjnych są schodkowe
  i nasycają się w 1 dla najbardziej skrajnych obserwacji testowych --
  dokładnie tam, gdzie PR-AUC jest najbardziej czuła.
  \item \textbf{Wagi}: różnice między schematami są małe (do 0,03), mniejsze
  niż wpływ wyboru detektorów. Przy fuzji wszystkich detektorów równe wagi
  są najgorsze, bo słaby detektor dostaje tyle samo co dobre.
  \item \textbf{Próg}: próg z walidacyjnego F1 daje średnio nieco wyższe F1
  (0,57--0,62) niż próg z udziału anomalii (0,56--0,61).
\end{itemize}

\subsection{Wagi i wybór detektorów w wariantach własnych}

\begin{figure}[H]
  \centering
  \includegraphics[width=\textwidth]{\fig{weights}}
  \caption{Średnie wagi detektorów (wszystkie podziały projektowe) w fuzji
  z wagami skill, w DSF i w CBF.}
  \label{fig:weights}
\end{figure}

Rysunek~\ref{fig:weights} pokazuje różnicę między wariantami. Wagi skill
rozkładają się na wszystkie detektory, także słabe (rekonstrukcja PCA).
DSF dzięki zachłannemu wyborowi zostawia zwykle jeden do trzech detektorów.
Na Cardio łączy OCSVM z innymi (w 8 z 10 podziałów), co daje najlepszy wynik
na tym zbiorze. Na WDBC zawsze zostawia jeden detektor (LOF w 18, IForest
w 12 podziałach), czyli wraca do problemu wyboru pojedynczego detektora na
2--3 anomaliach. CBF na Cardio we wszystkich 10 podziałach wybiera wyłącznie
OCSVM (przewaga na walidacji jest istotna), na WDBC uśrednia trzy lub cztery
detektory (bootstrap na 2--3 anomaliach nie odróżnia ich od najlepszego), na
Arrhythmii najczęściej IForest, LOF i OCSVM, a na Campaign IForest sam lub
z OCSVM.

\subsection{Przypadki poprawy i pogorszenia}
\label{sec:fail}

Tabela~\ref{tab:cases} porównuje decyzje zespołów z decyzjami najlepszego
pojedynczego detektora dla każdej obserwacji testowej.

\begin{table}[H]
  \centering\small
  \caption{Analiza przypadków (część testowa, suma po podziałach projektowych):
  anomalie wykryte przez zespół, a pominięte przez najlepszy pojedynczy
  detektor (,,uratowane''), i odwrotnie (,,utracone''), oraz fałszywe alarmy
  usunięte i dodane przez zespół.}
  \label{tab:cases}
  \resizebox{\textwidth}{!}{\tab{cases}}
\end{table}

\paragraph{Kiedy zespół pomaga.} Na Cardio DSF traci 37 anomalii, ratuje 23,
ale usuwa 57 z 88 fałszywych alarmów najlepszego detektora, dodając tylko
15. Precision rośnie z 0,81 do 0,89, a F1 z 0,79 do 0,81. Słabsze detektory
,,potwierdzają'' wysokie wyniki OCSVM tylko dla prawdziwych anomalii.
Na Arrhythmii CBF usuwa 33 z 89 fałszywych alarmów, dodając 13. Detektory są
tam porównywalnie dobre, a ich błędy w dużej mierze indywidualne, więc
uśrednienie zmniejsza wariancję wyniku -- klasyczny mechanizm działania
zespołów.

\paragraph{Kiedy zespół szkodzi.} Na WDBC fuzja średnia traci 15 z 79
anomalii, ratując 4. Mechanizm to \emph{rozmycie dobrego sygnału}: LOF
i OCSVM na pierwszej składowej PCA prawie bezbłędnie oddzielają guzy
złośliwe, a IForest i rekonstrukcja dodają do średniej szum, który przesuwa
część anomalii poniżej progu. Na Campaign MV 2/4 i fuzja średnia ratują
ok. 1400--1500 anomalii, ale dodają 4900--5800 fałszywych alarmów. Przy
11\,\% anomalii i słabych detektorach suma sygnałów obniża precyzję
(0,28--0,30) i fuzja średnia ma niższą PR-AUC niż najlepszy pojedynczy
detektor.

\paragraph{DSF a CBF.} Oba warianty własne mają niemal taką samą średnią
(PR-AUC 0,651 i 0,652), ale wygrywają w innych sytuacjach. DSF jest lepszy,
gdy walidacja jest duża, a słabsze detektory niosą informację uzupełniającą
(Cardio: 0,872 wobec 0,836; Campaign: 0,328 wobec 0,320). Zachłanny wybór
ma wtedy wiarygodną podstawę. CBF jest lepszy, gdy walidacja jest bardzo
mała (WDBC: 0,851 wobec 0,803). DSF dopasowuje tam skład zespołu do 2--3
anomalii i często wybiera gorszy detektor, a CBF w takiej sytuacji
uśrednia, bo nie ma podstaw, by odrzucić którykolwiek z dobrych
detektorów. Na Cardio konserwatyzm CBF jest wadą: istotna przewaga OCSVM
na walidacji sprawia, że odrzuca pozostałe detektory, choć w połączeniu
z OCSVM poprawiłyby wynik. Kryterium ,,czy detektor jest równie dobry
sam'' nie mierzy tego, czy jest użyteczny \emph{w zespole}.

\subsection{Test potwierdzający na nowych podziałach}
\label{sec:confirm}

\begin{table}[H]
  \centering\small
  \caption{Średnie PR-AUC i F1 (4 zbiory, część testowa) na podziałach
  projektowych (,,projekt.'') i na nowych podziałach (,,potw.'');
  ,,Lepszy niż najl.'' -- liczba zbiorów, na których PR-AUC wariantu jest
  wyższa niż najlepszego pojedynczego detektora.}
  \label{tab:confirm}
  \resizebox{\textwidth}{!}{\tab{confirmation}}
\end{table}

Na nowych podziałach (tabela~\ref{tab:confirm}) CBF ma najwyższą średnią
PR-AUC (0,628 wobec 0,620 najlepszego pojedynczego, 0,621 DSF i 0,620
wag skill). DSF i CBF przewyższają najlepszy pojedynczy detektor na 2 z 4
zbiorów i mają wyższe średnie F1 (0,580 i 0,579 wobec 0,565). Kierunek
efektów jest taki sam jak na podziałach projektowych: DSF wygrywa na Cardio
(0,851 wobec 0,801) i Campaign (0,334 wobec 0,309), CBF na WDBC (0,828
wobec 0,815) i Campaign (0,333). Fuzja średnia i MV 2/4 pozostają poniżej
najlepszego pojedynczego detektora w PR-AUC. Różnice średnich między
wariantami (0,01) są jednak mniejsze od odchyleń między podziałami na
Arrhythmii i WDBC (0,13--0,20). Wniosek ,,zespół jest średnio co najmniej
tak dobry jak najlepszy pojedynczy detektor'' jest potwierdzony, wniosek
,,jest lepszy'' -- tylko dla konkretnych zbiorów.

\subsection{Porównanie z zadaniami 1 i 2}

Tabela~\ref{tab:ladder} przedstawia PR-AUC uzyskaną na kolejnych etapach
projektu na tych samych podziałach. Każdy wybór dokonywany jest na
walidacji.

\begin{table}[H]
  \centering\small
  \caption{PR-AUC (test) na kolejnych etapach projektu: domyślny Isolation
  Forest w danych oryginalnych, wybór algorytmu w danych oryginalnych
  (zad.~2), wybór reprezentacji i algorytmu (zad.~2), najlepszy pojedynczy
  detektor bazowy oraz zespoły (zad.~3).}
  \label{tab:ladder}
  \resizebox{\textwidth}{!}{\tab{synthesis_ladder}}
\end{table}

\begin{figure}[H]
  \centering
  \includegraphics[width=\textwidth]{\fig{synthesis_ladder}}
  \caption{PR-AUC (test) na kolejnych etapach zadań 2 i 3.}
  \label{fig:ladder}
\end{figure}

Na zbiorach zadań 2 i 3 największy skok jakości daje \textbf{wybór
reprezentacji}: średnia PR-AUC rośnie z 0,544 do 0,663. Zysk pochodzi
głównie z Cardio (UMAP + OCSVM: 0,60 $\to$ 0,84) i WDBC (PCA: 0,71 $\to$
0,92). \textbf{Wybór algorytmu} w danych oryginalnych spośród trzech
detektorów nic tu nie daje (0,556 $\to$ 0,544). Isolation Forest jest
na tych zbiorach mocnym punktem wyjścia, a na WDBC walidacja wybiera
słabszy LOF. Zespół (zadanie~3) utrzymuje poziom najlepszego pojedynczego
detektora (CBF 0,652 wobec 0,654), a zwykła średnia go obniża (0,622).

\section{Synteza: co miało największy wpływ na jakość detekcji?}
\label{sec:synthesis}

Tabela~\ref{tab:effects} przedstawia dwa rodzaje miar wpływu na PR-AUC.
Pierwszy to \emph{rozrzut} (maksimum $-$ minimum), który mówi, ile można
zyskać lub stracić danym wyborem. Drugi to \emph{zrealizowany zysk}
z zespołu przy wyborze na walidacji.

\begin{table}[H]
  \centering\small
  \caption{Wpływ czynników na PR-AUC (test): rozrzut między 3 detektorami
  w danych oryginalnych (zad.~2), średni rozrzut jednego detektora między
  reprezentacjami (zad.~2) oraz zysk zespołu względem najlepszego
  pojedynczego detektora wybranego na walidacji (zad.~3).}
  \label{tab:effects}
  \resizebox{0.85\textwidth}{!}{\tab{synthesis_effects}}
\end{table}

Wybór algorytmu badano też w zadaniu~1 na pięciu innych zbiorach
z szerszą siatką pięciu metod. Tam rozrzut PR-AUC między metodami był
największy ze wszystkich czynników: średnio 0,40, na Breastw 0,64, na Wine
0,61, na Cardio 0,56. Zrealizowany zysk z wyboru metody na walidacji
względem domyślnego Isolation Forest wyniósł średnio +0,15 (0,511 $\to$
0,658), w tym Wine +0,51 (OCSVM) i Cardio +0,16 (OCSVM).

\paragraph{Odpowiedź.} Na podstawie wyników trzech zadań:
\begin{enumerate}
  \item \textbf{Wybór algorytmu i reprezentacji mają porównywalny,
  największy wpływ -- ale nie da się ich rozdzielić.} W zadaniu~1 wybór
  metody podniósł średnią PR-AUC o 0,15, a w zadaniu~2 wybór reprezentacji
  (łącznie z algorytmem w niej) o 0,12. Oba efekty są skupione na kilku
  zbiorach (Wine, Cardio, WDBC), a w obu przypadkach kluczowa jest
  \emph{para}: ta sama reprezentacja pomaga jednemu detektorowi i szkodzi
  innemu (UMAP: OCSVM na Cardio +0,24, LOF $-0{,}07$), a ten sam algorytm
  zmienia ranking w zależności od reprezentacji. Detektor to w praktyce
  para reprezentacja--algorytm.
  \item \textbf{Reprezentacja ma duży rozrzut (średnio 0,19), w obie strony.}
  PCA z jedną składową daje na WDBC +0,21 dla LOF, UMAP z OCSVM na Cardio
  +0,24, a UMAP z IForest na Arrhythmii $-0{,}26$. Wybór na walidacji
  przenosi się na test tam, gdzie walidacja ma dość anomalii.
  \item \textbf{Połączenie detektorów ma najmniejszy wpływ.} Najlepsze
  warianty (CBF, DSF) zmieniają średnią PR-AUC o ok. 0,00 (na nowych
  podziałach $+0{,}01$), a zwykła średnia o $-0{,}03$. Zespół pomaga, gdy
  detektory wnoszą różne informacje, a walidacja pozwala to ocenić (Cardio,
  Arrhythmia, Campaign: +0,01 do +0,04), i szkodzi, gdy dobry sygnał
  zostaje rozmyty słabszymi detektorami (WDBC: do $-0{,}10$).
  \item Ponad wszystkimi trzema czynnikami stoi \textbf{natura danych}:
  rozrzut PR-AUC między zbiorami (0,33--0,93 dla najlepszej metody) jest
  większy niż efekt któregokolwiek wyboru. Na Campaign, gdzie ,,anomalie''
  (klienci zakładający lokatę) nie są wyjątkami statystycznymi, ani
  algorytm, ani reprezentacja, ani zespół nie podniosły PR-AUC powyżej 0,34
  przy bazowym 0,11.
\end{enumerate}

\section{Rekomendacje: kiedy zespół, a kiedy pojedynczy detektor}

\begin{itemize}
  \item \textbf{Pojedynczy detektor}, gdy jeden detektor wyraźnie dominuje
  i ma unikalną wiedzę, a pozostałe są słabe i z nim skorelowane, albo gdy
  walidacja jest zbyt mała, by ocenić detektory (WDBC). Uśrednianie rozmyje
  wtedy dobry sygnał.
  \item \textbf{Fuzja z zachłannym wyborem (DSF, wagi skill)}, gdy walidacja
  ma dziesiątki anomalii, a detektory działają w różnych reprezentacjach
  i wykrywają różne anomalie (Cardio, Campaign).
  \item \textbf{CBF}, gdy walidacja jest mała lub nie wiadomo, czy zespół
  pomoże. Z konstrukcji sprowadza się do najlepszego detektora, gdy ten
  wyraźnie dominuje, i do średniej, gdy detektory są nieodróżnialne. Na nowych
  podziałach miał najwyższą średnią PR-AUC.
  \item \textbf{Majority Voting z progiem 2/4 lub 3/4}, gdy priorytetem jest
  stabilna decyzja binarna, a nie ranking (najwyższe F1 na Campaign).
  \item \textbf{Nie stosować} fuzji przez maksimum ani średniej wszystkich
  detektorów bez kontroli jakości. Normalizację dopasowywać na wynikach
  \emph{out-of-sample} (walidacja), a nie na wynikach części uczącej.
\end{itemize}

\section{Ograniczenia}

\begin{itemize}
  \item Części walidacyjne WDBC zawierają 2--3 anomalie, co ogranicza
  wiarygodność każdej metody doboru wag i progów, a nawet wyboru
  pojedynczego detektora. Trzydzieści podziałów zmniejsza niepewność
  średniej, ale nie poprawia decyzji w pojedynczym podziale.
  \item Dla Campaign użyto 5 podziałów projektowych i 5 potwierdzających
  (koszt obliczeń). Odchylenia między podziałami są tam małe (0,01--0,03).
  \item Detektory bazowe przejęto z zadania~2 razem z jego decyzjami. Na
  WDBC LOF i OCSVM są niemal identyczne, więc zespół ma de facto trzy różne
  detektory.
  \item Zbadano zespoły heterogeniczne złożone z 4 detektorów. Nie badano
  zespołów jednorodnych (np. wielu LOF o różnych $k$) ani uczenia
  meta-modelu (stacking) na walidacji, który przy małych walidacjach
  byłby silnie przeuczony.
\end{itemize}

\section{Wnioski}

\begin{enumerate}
  \item Zespół detektorów nie jest automatycznie lepszy od najlepszego
  detektora. Najlepsze warianty (DSF, CBF) są średnio tak samo dobre jak
  najlepszy pojedynczy detektor wybrany na walidacji. Wygrywają na Cardio,
  Arrhythmii i Campaign (o 0,01--0,04 PR-AUC), a przegrywają na WDBC.
  \item O efekcie zespołu decyduje relacja między jakością
  a komplementarnością detektorów oraz wielkość walidacji. Komplementarność
  pomaga, gdy słabsze detektory potwierdzają trafienia dobrego, a szkodzi,
  gdy rozmywają jego prawie bezbłędny sygnał (WDBC).
  \item Wybór detektorów do zespołu jest ważniejszy niż schemat wag,
  a normalizację trzeba dopasować na wynikach out-of-sample. Wyniki
  in-sample są przesunięte (Campaign: 71\,\% wyników testowych OCSVM powyżej
  maksimum wyników uczących), a fuzja przez maksimum z rangami dopasowanymi
  na nich traci 0,11 PR-AUC.
  \item Własne warianty DSF i CBF realizują przeciwne strategie (zaufanie
  walidacji i ostrożność wobec niej) i sprawdzają się w różnych warunkach:
  DSF przy dużej walidacji, CBF przy małej. CBF miał najwyższą średnią
  PR-AUC w teście potwierdzającym.
  \item W skali całego projektu o jakości detekcji decydowały kolejno:
  natura danych, dobór pary algorytm--reprezentacja i dopiero na końcu
  połączenie detektorów.
\end{enumerate}

\section{Odtworzenie wyników}

\begin{verbatim}
python -m task3_ensemble.run_base                    # wymaga wyników zad. 2
python -m task3_ensemble.run_ensembles
python -m task3_ensemble.run_base --confirm          # nowe ziarna
python -m task3_ensemble.run_ensembles --confirm
python -m task3_ensemble.analyze
\end{verbatim}

% Sources of the benchmark datasets, shared by the three reports.
\begin{thebibliography}{99}
\small

\bibitem{adbench} S.~Han, X.~Hu, H.~Huang, M.~Jiang, Y.~Zhao.
\emph{ADBench: Anomaly Detection Benchmark}. Advances in Neural Information
Processing Systems 35 (NeurIPS 2022), Datasets and Benchmarks Track, 2022.
Repozytorium \url{https://github.com/Minqi824/ADBench},
commit \code{3dac822}.

\bibitem{odds} S.~Rayana. \emph{ODDS Library}. Stony Brook University,
Department of Computer Science, 2016. \url{https://odds.cs.stonybrook.edu}.

\bibitem{dami} G.~O.~Campos, A.~Zimek, J.~Sander, R.~J.~G.~B.~Campello,
B.~Micenková, E.~Schubert, I.~Assent, M.~E.~Houle. \emph{On the evaluation of
unsupervised outlier detection: measures, datasets, and an empirical study}.
Data Mining and Knowledge Discovery 30(4), 891--927, 2016.
\url{https://www.dbs.ifi.lmu.de/research/outlier-evaluation/}.

\bibitem{uci} M.~Kelly, R.~Longjohn, K.~Nottingham. \emph{The UCI Machine
Learning Repository}. University of California, Irvine.
\url{https://archive.ics.uci.edu}.

\bibitem{breastw} W.~H.~Wolberg, O.~L.~Mangasarian. \emph{Multisurface method
of pattern separation for medical diagnosis applied to breast cytology}.
Proceedings of the National Academy of Sciences 87(23), 9193--9196, 1990.
(UCI: Breast Cancer Wisconsin (Original)).

\bibitem{wdbc} W.~N.~Street, W.~H.~Wolberg, O.~L.~Mangasarian. \emph{Nuclear
feature extraction for breast tumor diagnosis}. IS\&T/SPIE Symposium on
Electronic Imaging: Science and Technology, vol.~1905, 861--870, 1993.
(UCI: Breast Cancer Wisconsin (Diagnostic)).

\bibitem{mammography} K.~S.~Woods, C.~C.~Doss, K.~W.~Bowyer, J.~L.~Solka,
C.~E.~Priebe, W.~P.~Kegelmeyer. \emph{Comparative evaluation of pattern
recognition techniques for detection of microcalcifications in mammography}.
International Journal of Pattern Recognition and Artificial Intelligence
7(6), 1417--1436, 1993.

\bibitem{wine} S.~Aeberhard, M.~Forina. \emph{Wine}. UCI Machine Learning
Repository, 1991.

\bibitem{cardio} D.~Ayres-de Campos, J.~Bernardes, A.~Garrido,
J.~Marques-de-Sá, L.~Pereira-Leite. \emph{SisPorto 2.0: a program for
automated analysis of cardiotocograms}. Journal of Maternal-Fetal Medicine
9(5), 311--318, 2000. (UCI: Cardiotocography).

\bibitem{campaign} S.~Moro, P.~Cortez, P.~Rita. \emph{A data-driven approach
to predict the success of bank telemarketing}. Decision Support Systems 62,
22--31, 2014. (UCI: Bank Marketing).

\bibitem{arrhythmia} H.~A.~Güvenir, B.~Acar, G.~Demiröz, A.~Çekin.
\emph{A supervised machine learning algorithm for arrhythmia analysis}.
Computers in Cardiology 1997, 433--436. (UCI: Arrhythmia).

\bibitem{heart} \emph{Statlog (Heart)}. UCI Machine Learning Repository.

\end{thebibliography}


\end{document}
